\subsection{Motivación y contexto de selección}

Un simulador tipo Schrödinger puede
particionar el vector de estado en bloques independientes, mantener una caché de bloques descomprimidos y
aplicar una política de descompresión bajo demanda para evitar reconstruir el estado completo en cada operación.
Sin embargo, para materializar esa arquitectura en un simulador real es necesario seleccionar una biblioteca
de compresión concreta que se adapte bien al acceso parcial por bloques y al patrón repetido de lectura,
modificación y escritura de amplitudes.

Por esta razón, la decisión relevante es qué biblioteca de compresión sin pérdida facilita 
una integración eficiente en una ruta crítica de simulación cuántica. En este contexto, interesan
especialmente las soluciones que trabajen con bloques independientes, permitan incorporar filtros previos sobre
datos de punto flotante y minimicen la latencia de descompresión cuando solo una fracción del vector de estado
debe pasar temporalmente a representación densa.